\begin{lemma}
	\label{lemma:main-ec}
	If $\pker$ is a lossy and \wlap{} public-key encryption scheme, $\pkem$ is a lossy and $\mathcal{U}_{\msgspace_\mat}$-\slap{} public-key encryption scheme with message space $\msgspace_\mat$, $(\csetup, \commit)$ satisfies statistical hiding and computational binding, $\Pi_{E}$ satisfies zero-knowledge, and $\Pi_{KG}$ satisfies strong simulation extractability, then \schemename{} satisfies escrow hiding.
\end{lemma}

\begin{proof}
	We start by describing our simulator $\Sim$, and use use a sequence of games to show that:
	$$\bigg\{ \escrow(\ppbpp, pk, x, r) \bigg\} \cindist \bigg\{\Sim(\ppbpp, \td, pk, f(t, x)) \bigg\}$$

	Our simulator behaves as follows:

	\begin{pchstack}[center, space=1em]
		\procedure[linenumbering]{$\Sim(\ppbpp, \td, pk, z = f(t, x))$}{
			\pcparse (\secpar, \compp, \pkerpp, \pkempp, \crs_{KG}, \crs_{E}, n) = \ppbpp \\
			\pcif z = (\msg, \uval') \:\pcreturn \escrow(\ppbpp, pk, (\uval'||0^{n - |\uval'|}, \msg)) \\
			\pcelse \pcif z = \bot \:\pcreturn \escrow(\ppbpp, pk, (0, \epsilon))
		}
	\end{pchstack}

	We begin our analysis with the case that $z = \bot$.
	Additionally, as $z = \bot$, it is the case that $\uval \leq \thr$.
	Therefore, there either exists some index $d$ such that $\forall i < d$, $\uval[i] = \thr[i]$, and $\uval[d] < \thr[d]$, or $\uval = \thr$.
	We first prove security in the case that $\uval < t$ via the following sequence of hybrids.

	\begin{itemize}
		\item $\newhyb$: The original escrow hiding game when $b = 0$, i.e. $\escrow(\ppbpp, pk, x, r)$.

		\item $\nexthyb$: The proof $\zkproof_{E}$ is simulated using the trapdoor $\td_{E}$. By the zero-knowledge property of $\Pi_{E}$, $\prevhyb\pindist \currhyb$.

		\item $\nexthyb$: Replace $\com_{E}$ with $\commit(\compp, (0, \epsilon))$. By the statistical hiding property of the commitment scheme scheme, $\prevhyb \sindist \currhyb$.

		\item $\nexthyb$: Replace $\otpct_d^{\mat}$ with an element sampled uniformly at random from $\msgspace_\mat$.
		      Let $\nu(\secpar)$ be an adversary's advantage in distinguishing between $\prevhyb$ and $\currhyb$.
		      We will show that $\nu(\secpar)$ must be negligible for all adversaries $\adv$.

		      The adversary $\adv$ distinguishing between the two hybrids outputs $(\thr, r_{KG}, pk = (pk', C_{KG}, \pi_{KG}), x)$ in its first stage.
		      We can divide our analysis into the following three cases depending on the properties of these outputs.
		      \begin{enumerate}
			      \item There exist some $r_{KG}'$ and $\thr' \neq \thr$ such that \\$\commit(\thr, r_{KG}) = \commit(\thr', r_{KG}')$ and $(\cdot, pk') = \kgen'(\secpar, \pkepp^{\rev}, \pkepp^{\mat}, \numdigits, \thr')$.
			      \item $(\cdot, pk') \neq \kgen'(\secpar, \pkepp^{\rev}, \pkepp^{\mat}, \numdigits, \thr)$ and (1) does not hold.
			      \item Neither (1) nor (2) holds.
		      \end{enumerate}

		      Let the probability of the outputs satisfying property $j$ be denoted by $\nu_j(\secpar)$.
		      If $\nu(\secpar)$ is non-negligible, then at least one of $\nu_1(\secpar)$, $\nu_2(\secpar)$, $\nu_3(\secpar)$ must be non-negligible.

		      ${}$

		      Suppose that $\nu_2(\secpar)$ is non-negligible.
		      In this case $\adv$ produced a proof of a false statement and we can construct a reduction $\bdv$ to the simulation-extractability of $\Pi_{KG}$.
		      $\bdv$ runs $\adv$, receives $(\thr, r_{KG}, pk = (pk', C_{KG}, \pi_{KG}), x)$, and outputs $((pk', C_{KG}), \pi_{KG})$.
		      As the statement is false the extractor must fail, meaning $\prob{\bdv \text{ wins}}$ is non-negligible, thus contradicting our assumption that $\Pi_{KG}$ satisfies strong simulation-extractability.
		      Therefore, $\nu_2(\secpar)$ is negligible.

		      ${}$

		      As $\nu_2(\secpar)$ is negligible, we can assume that there exists an extractor $\Ext$ that given $\pi_{KG}$ outputs $(\thr', sk', r_k, r_c)$ such that $C_{KG} = \commit(\compp, \thr; r_c)$ and $(sk', pk') = \kgen'(\secpar, \pkepp^{\rev}, \pkepp^{\mat}, \base, \numdigits, \thr; r_k)$.

		      ${}$

		      Suppose $\nu_1(\secpar)$ is non-negligible.
		      In this case we can break the computational binding property of the commitment scheme by using a reduction that uses the extractor and outputs $(t, r_{KG}, t', r_c)$.
		      Therefore, $\nu_1(\secpar)$ must be negligible.

		      ${}$

		      Suppose $\nu_3(\secpar)$ is non-negligible.
		      In this case we can construct a reduction $\bdv$ to the computational lossiness of $\pkem$.
		      $\bdv$ runs $\adv$, receives $(\thr, r_{KG}, pk = (pk', C_{KG}, \pi_{KG}), x)$, and uses $\Ext$ to extract $(\thr, sk', r_k, r_c)$.
		      $\bdv$ then uses $r_k$ to infer $r$ such that $(\cdot, \matpk_d^{\uval[d]}) = \pkem.\kgenlossy(\matpp; r)$ (the key must be lossy as $\uval[d] < \thr[d]$),
		      samples $\otpct'$ uniformly at random from $\msgspace_\mat$, and sends $(r, \otpct_d^{\mat}, \otpct')$ to its challenger.
		      $\bdv$ receives $\ct$, sends to $\adv$ an escrow value identical to the one produced in $\currhyb$ with the exception that $\mct_{d}$ is replaced by $\ct$, and outputs what $\adv$ outputs.

		      In the case that $\bdv$ is in $\complossygame_0$, $\adv$'s view is identical to game $\prevhyb$, and in $\complossygame_1$ $\adv$'s view is identical to $\currhyb$.
		      We therefore conclude that $\bdv$'s advantage is equal to $\nu_3(\secpar)$, and so $\nu_3(\secpar)$ must be negligible.

		      ${}$

		      As $\nu_1(\secpar)$, $\nu_2(\secpar)$, and $\nu_3(\secpar)$ are all negligible, $\nu(\secpar)$ must be negligible.
		      Therefore, we have that $\prevhyb \cindist \currhyb$.

		      Identical logic showing that $\nu_0(\secpar)$ and $\nu_1(\secpar)$ must be negligible will be implicitly used in future arguments that depend upon properties of $pk'$.
		      For brevity we will assume the existence of an extractor outputting $(t, sk', r_k, r_c)$ and focus only on $\nu_3(\secpar)$ without repeating the full proof.

		      \newcounter{DBM}
		      \setcounter{DBM}{\value{hybctr}}

		\item $\nexthyb$: In this hybrid $\otpct_i^{\mat}$ is replaced by an element sampled uniformly at random from $\msgspace_\mat$ for all $i \in [d + 1, \numdigits - 1]$.
		      We show that $\prevhyb \pindist \currhyb$ via a series of sub-hybrids.
		      In hybrid $\hyb'_k$ for $k \in [d + 1, \numdigits - 1]$, replace $\otpct_k^{\mat}$ as described.
		      Clearly $\hyb'_{d}$ is equivalent to $\prevhyb$ and $\hyb'_{\numdigits - 1}$ is equivalent to $\currhyb$.
		      As the output contains no information regarding $\otpk_{k - 1}$ (as $\otpct^{\mat}_{k - 1}$ has been replaced with a random element in the previous hybrid), $\hyb'_{k} \cindist \hyb'_{k + 1}$ for all $k$, and so $\prevhyb \pindist \currhyb$.
		      \newcounter{HighBM}
		      \setcounter{HighBM}{\value{hybctr}}

		\item $\nexthyb$: In this hybrid $\otpct_i^{\mat}$ is replaced by an element sampled uniformly at random from $\msgspace_\mat$ for all $i \in [d - 1, 1]$.
		      We show that $\prevhyb \pindist \currhyb$ via a series of sub-hybrids.
		      In hybrid $\hyb'_k$ for $k \in [d - 1, 1]$, replace $\otpct_k^{\mat}$ as described.

		      In $\hyb'_k$ we replace  $\otpk_{k-1} \cdot \otpk_{k}$ with a random element.
		      Then, as the output contains no information regarding $\otpk_{k}$ (as $\otpct^{\mat}_{k + 1}$ has been replaced with a random element in the previous hybrid, or in $\hyb_3$ for $k = d - 1$), $\hyb'_k \pindist \hyb'_{k - 1}$ for all $k$.

		      ${}$

		      Note that at this point all all values $\otpct^{\mat}_i$ have been replaced with uniformly random elements of $\msgspace_\mat$.
		      \newcounter{LowBM}
		      \setcounter{LowBM}{\value{hybctr}}

		\item $\nexthyb$: In this hybrid $\mct_i$ is replaced by $\pkem.\enc(\pkempp, \matpk_i^{0}, \otpct_i^{\mathsf{M}})$ for all $i \in [d - 1]$.
		      We show that $\prevhyb \cindist \currhyb$ via a series of sub-hybrids.
		      In sub-hybrid $\hyb'_k$ for $k \in [d-1]$, replace $\mct_k$ as described.
		      In the case that $\thr[k] = \uval[k] = 0$, $\hyb'_k$ is identical to the previous one.
		      In the case that $\thr[k] = \uval[k] > 0$, note that $\matpk_i^{v[i]}$ is injective and $\matpk_i^0$ is lossy.
		      Skipping the logic gone over in $\hyb_3$, we will now show that in this case $\nu_3(\secpar)$ is negligible.
		      To do so, we will construct a reduction $\bdv$ to the strong lossy ambiguity property of $\pkem$.

		      ${}$

		      $\bdv$ runs $\adv$, receives $(\thr, r_{KG}, pk = (pk', C_{KG}, \pi_{KG}), x)$, and uses $\Ext$ to extract $(\thr, sk', r_k, r_c)$.
		      $\bdv$ then uses $r_k$ to infer $r_0$, $r_{\uval}$ such that $(\cdot, \matpk_k^{0}) \allowbreak=\allowbreak \pkem.\kgenlossy\allowbreak(\pkepp^{\mat},\allowbreak \secparam; r_0)$ and $(\cdot, \matpk_k^{\uval[i]}) = \pkem.\kgen(\pkepp^{\mat}, \secparam; r_{\uval})$, and sends $(r_{\uval}, r_0)$ to its challenger.
		      $\bdv$ receives $\ct^*$, sends to $\adv$ an escrow value identical to the one produced in
		      $\hyb'_{k}$ with the exception that $\mct_{k}$ is replaced by $\ct^*$, and outputs what $\adv$ outputs.
		      In the case that $\bdv$ is in $\slambgame_0$, $\adv$'s view is identical to sub-hybrid $\hyb'_{k-1}$, and in $\slambgame_1$ $\adv$'s view is identical to sub-hybrid $\hyb'_{k}$.
		      We therefore conclude that $\bdv$'s advantage is equal to $\nu_3(\secpar)$, and so $\nu_3(\secpar)$ must be negligible.
		      Therefore, we have that $\hyb'_k \cindist \hyb'_{k + 1}$ for all $k$, and so $\prevhyb \cindist \currhyb$.

		      Future transitions that rely on the weak/strong lossy ambiguity of \\$\pker$/$\pkem$ will proceed in the same way, and so we will omit the full proofs for brevity.
		\item $\nexthyb$: In this hybrid we replace $\mct_d$ with $\pkem.\enc(\pkempp, \matpk_d^{0}, \otpct_d^{\mathsf{M}})$.
		      In the case that $\uval[d] = 0$, this hybric is identical to $\prevhyb$.
		      Otherwise, as $\uval[d] < \thr[d]$, we have that $\matpk_d^{\uval[d] + 1}$ is lossy.
		      As such, by the weak lossy ambiguity property of $\pkem$, $\prevhyb \cindist \currhyb$.

		\item $\nexthyb$: In this hybrid we replace $\mct_i$ with $\pkem.\enc(\pkepp, \matpk_i^0, \otpct_i^{\mat r})$ for all $i \in [d + 1, \numdigits - 1]$.
		      We show that $\prevhyb \cindist \currhyb$ via a series of sub-hybrids.
		      In sub-hybrid $\hyb'_k$ for $k \in [d + 1, \numdigits - 1]$, replace $\mct_k$ as described.
		      % Clearly $\hyb'_d$ is equivalent to $\prevhyb$ and $\hyb'_{\numdigits - 1}$ is equivalent to $\currhyb$.

		      In the case that $\uval[k] = 0$, the sub-hybrid is identical to the previous sub-hybrid.
		      Otherwise, we can split the possibilities into several cases.
		      \begin{enumerate}
			      \item When $\uval[k] \geq \thr[k] > 0$, we have that $\matpk_k^{\uval[k] + 1}$ is injective while $\matpk_k^{0}$ is lossy.
			            Indistinguishability therefore follows from the strong lossy ambiguity property of $\pkem$.

			      \item When $\thr[k] > \uval[k] > 0$ we have that $\matpk_k^{\uval[k] +1}$ and $\matpk_k^0$ are both lossy. Indistinguishability therefore follows from the weak lossy ambiguity property of $\pkem$.
		      \end{enumerate}
		      We therefore have that $\prevhyb \cindist \currhyb$ by the strong lossy ambiguity of $\pkem$.

		      Note that at this point $\mct_i$ is encrypted under $\matpk_{i}^0$ for all $i$.

		\item $\nexthyb$: Replace $\rct_i$ with $\pker.\enc(\pkerpp, \revpk_i^0, \otpk^*)$ for all $i \in [d]$.
		      We show that $\prevhyb \cindist \currhyb$ via a series of sub-hybrids.
		      In sub-hybrid $\hyb'_k$ for $k \in [d]$, replace $\rct_k$ as described.
		      In the case that $\uval[k] = 0$ the sub-hybrid is identical to the previous sub-hybrid.
		      Otherwise, as $\thr[k] \geq \uval [k]$ we have that $\revpk_k^{\uval[k]}$ is lossy, and so $\hyb'_k \cindist \hyb'_{k + 1}$ for all $k$ by the weak lossy ambiguity of $\pker$.

		\item $\nexthyb$: Replace $\otpct_i^\rev$ with an element sampled uniformly at random from $\msgspace_\mat$ for all $i \in [d + 1, \numdigits]$.
		      We can show that $\prevhyb \pindist \currhyb$ via a series of sub-hybrids, where at sub-hybrid $\hyb'_k$ we replace $\otpct_k^{\rev}$ as described.
		      As the output contains no information regarding $\otpk_{k - 1}$ (as $\otpct_{k - 1}^{\mat}$ has been replaced with a random element), $\hyb'_{k} \pindist \hyb'_{k+1}$.
		      Therefore $\prevhyb \pindist \currhyb$.
		      \newcounter{RevOTPCtDtoN}
		      \setcounter{RevOTPCtDtoN}{\value{hybctr}}

		\item $\nexthyb$: Replace $\rct_i$ with $\pker.\enc(\pkerpp, \revpk_i^0, \epsilon)$ for all $i \in [d + 1, \numdigits]$. As the output contains no information regarding $\rct_i$ (as $\otpct^\rev_i$ was previously replaced with a random element), we have that $\prevhyb \pindist \currhyb$.

		\item $\nexthyb$: Replace $\rct_i$ with $\pker.\enc(\pkerpp, \revpk_i^0, \epsilon)$ for all $i \in [d]$.
		      We can show that $\prevhyb \cindist \currhyb$ via a series of sub-hybrids, where at sub-hybrid $\hyb'_k$ we replace $\rct_k$ as described.
		      As $\revpk_k^0$ is always lossy, and $\revpk_k^{v[k]}$ is lossy by our assumption that $v \leq \thr$, by the computational lossiness property of $\pker$ we have that $\hyb'_{k} \cindist \hyb'_{k + 1}$ for all $k$, and so $\prevhyb \cindist \currhyb$.

		      Note that at this point $\rct_i$ encrypts $\epsilon$ under $\revpk_i^0$ for all $i$.

		\item $\nexthyb$: Replace $\otpct^*$ with $\otpk^* \cdot \epsilon$.
		      As the output contains no information about $\otpk^*$ (as all $\rct_i$ now encrypt $\epsilon$), we have that $\prevhyb \pindist \currhyb$.

		\item $\nexthyb$: This hybrid is identical to the functioning of the simulator.
		      We have that $\prevhyb \cindist \currhyb$ via a series of hybrids that reverts the changes made in hybrids so far as necessary.
		      Indistinguishability follows from the same justifications as above.
	\end{itemize}

	We therefore conclude that $$\bigg\{ \escrow(\ppbpp, pk, (\uval, \msg), r) \bigg\} \cindist \bigg\{\Sim(\ppbpp, \td, pk, f(\thr, (\uval, \msg)) \bigg\}$$ in the case that $\uval < \thr$.

	We now turn to the case that $\uval = \thr$, and consider the following series of hybrids.
	We keep things here at the proof-sketch level and refer the reader to the previous series of hybrids for more detailed reductions.

	\begin{itemize}
		\item $\newhyb$: The original escrow hiding game when $b = 0$, i.e. $\escrow(\ppbpp, pk, x, r)$.

		\item $\nexthyb$: The proof $\zkproof_{E}$ is simulated using the trapdoor $\td_{E}$. By the zero-knowledge property of $\Pi_{E}$, $\prevhyb \pindist \currhyb$.

		\item $\nexthyb$: Replace $\com_{E}$ with $\commit(\compp, (0, \epsilon))$. By the statistical hiding property of the commitment scheme scheme, $\prevhyb \sindist \currhyb$.

		\item $\nexthyb$: Replace $\revct_i$ with $\pker.\enc(\pkerpp, \revpk_i^0, \otpk^*)$ for all $i \in [\numdigits]$.
		      As $\revpk_i^{\uval[i]}$ is lossy for all $i$ (as $\uval = \thr$), we have that $\prevhyb \cindist \currhyb$ via a series of sub-hybrids and the \wla{} of $\pker$.

		\item $\nexthyb$: Replace $\revct_i$ with $\pker.\enc(\pkerpp, \revpk_i^0, \epsilon)$ for all $i \in \numdigits$.
		      As $\revpk_i^0$ is lossy for all $i$, we have that $\prevhyb \cindist \currhyb$ via a series of sub-hybrids and the computational lossiness of $\pker$.

		\item $\nexthyb$: Replace $\otpct^*$ with $\otpk^* \cdot \epsilon$.
		      As the output contains no information about $\otpk^*$, we have that $\prevhyb \pindist \currhyb$ by the perfect security of $\otp'$.

		\item $\nexthyb$: Replace $\mct_i$ with $\pkem.\enc(\pkempp, \matpk_i^0, \otpct_i^{\mat})$ for all $i \in [\numdigits]$.
		      As $\matpk_i^{\uval[i] + 1}$ is injective for all $i$ (as $\uval = \thr$), we have that $\prevhyb \cindist \currhyb$ via a series of sub-hybrids and the \sla{} of $\pkem$.

		\item $\nexthyb$: The simulator honestly generates the proof $\pi_{E}$.
		      By the zero-knowledge property of $\Pi_{E}$, $\prevhyb \pindist \currhyb$.
	\end{itemize}
	Note that $\currhyb$ is identical to the functioning of the simulator.


	Finally, we consider the case that $\uval > \thr$.
	If $\uval[i] \geq \thr[i]$ for all $i \in [\numdigits]$ then $\uval' = \uval$, and so the simulation is clearly perfect.
	Otherwise, there must be some digit $d$ for which $\uval[d] < \thr$.
	We can then use an identical sequence of games as in the case that $\uval < \thr$ to show that the output of the simulator is indistinguishable to the output of $\escrow$.

	We therefore conclude that $$\bigg\{ \escrow(\ppbpp, pk, (\uval, \msg), r) \bigg\} \cindist \bigg\{\Sim(\ppbpp, \td, pk, f(\thr, (\uval, \msg)) \bigg\}$$ for $\uval < \thr$, $\uval = \thr$, and $\uval > \thr$, and therefore \schemename{} satisfies escrow hiding.
\end{proof}

\cref{thm:ppbs-main} then follows directly from \cref{lemma:main-soundness}, \cref{lemma:main-correctness}, \cref{lemma:main-th}, and \cref{lemma:main-ec}.
